\documentclass[12pt,a4paper]{article}
\usepackage{amsmath, amssymb}
\usepackage[margin=2.5cm]{geometry}
\usepackage{enumitem}
\usepackage{xcolor}

\newcommand{\easy}{\textcolor{green!60!black}{$\bigstar$}}
\newcommand{\medium}{\textcolor{orange}{$\bigstar\bigstar$}}
\newcommand{\hard}{\textcolor{red}{$\bigstar\bigstar\bigstar$}}

\title{\textbf{Solutions: Additional Exercises 3}\\Power Series}
\author{Series and Integral Transforms}
\date{}

\begin{document}
\maketitle

% -------------------------------------------------------
\section*{Exercise 1 — Radius and Domain of Convergence}

\textit{Method}: find radius $R$ via Ratio or Root test, then check each endpoint separately.

\subsection*{(a) \easy\ $\displaystyle\sum_{n=1}^{\infty}\frac{x^n}{3^n \cdot n}$}

Coefficients $a_n=1/(3^n n)$.  By the Ratio Test:
$R = \lim\dfrac{|a_n|}{|a_{n+1}|} = \lim\dfrac{3^{n+1}(n+1)}{3^n n} = 3$.

\textit{Endpoints}:
\begin{itemize}
  \item $x=3$: $\sum 1/n$ (harmonic) — \textbf{diverges}.
  \item $x=-3$: $\sum(-1)^n/n$ — Leibniz, \textbf{converges}.
\end{itemize}
\textbf{Domain}: $[-3, 3)$.

\subsection*{(b) \easy\ $\displaystyle\sum_{n=0}^{\infty}\frac{x^n}{n!}$}

Ratio Test: $\lim\dfrac{n!}{(n+1)!}|x|=\lim\dfrac{|x|}{n+1}=0$ for all $x$.  So $R=+\infty$.

\textbf{Domain}: $(-\infty, +\infty)$.

\subsection*{(c) \easy\ $\displaystyle\sum_{n=1}^{\infty}\frac{(-1)^n x^n}{n}$}

$R=1$ (Ratio Test).

\textit{Endpoints}:
\begin{itemize}
  \item $x=1$: $\sum(-1)^n/n$ — Leibniz, \textbf{converges}.
  \item $x=-1$: $\sum(-1)^n(-1)^n/n=\sum 1/n$ — \textbf{diverges}.
\end{itemize}
\textbf{Domain}: $(-1, 1]$.

\subsection*{(d) \easy\ $\displaystyle\sum_{n=0}^{\infty} n!\, x^n$}

Ratio Test: $\lim(n+1)|x|=+\infty$ for $x\neq 0$.  So $R=0$.

\textbf{Domain}: $\{0\}$.

\subsection*{(e) \medium\ $\displaystyle\sum_{n=1}^{\infty}\frac{(-1)^n x^{2n-1}}{5^n\sqrt{n+1}}$}

The series has only odd powers of $x$.  Write $u = x^2$.  The coefficients (of $u^n$) satisfy:
\[
  R_u = \lim\frac{5^{n+1}\sqrt{n+2}}{5^n\sqrt{n+1}} = 5.
\]
So $|x^2| < 5$, meaning $|x| < \sqrt{5}$.

\textit{Endpoints} $x=\pm\sqrt{5}$: the series becomes $\sum \pm 1/\sqrt{n+1}$, which diverges ($p=1/2$).

\textbf{Domain}: $(-\sqrt{5},\,\sqrt{5})$.

\subsection*{(f) \medium\ $\displaystyle\sum_{n=0}^{\infty}\frac{2^n x^n}{(2n+1)^2 n^{1/4}}$}

Ratio Test gives $R=1/2$.

\textit{Endpoints} $x=\pm 1/2$: the series becomes $\sum\dfrac{\pm 1}{(2n+1)^2 n^{1/4}}$. Since $(2n+1)^2 n^{1/4}\sim 4n^{9/4}$ and $p=9/4>1$, this \textbf{converges absolutely}.

\textbf{Domain}: $[-\tfrac{1}{2}, \tfrac{1}{2}]$.

\subsection*{(g) \medium\ $\displaystyle\sum_{n=1}^{\infty}\frac{(-1)^{n-1} x^{2n-2}}{3^{n-1}\cdot n\sqrt{n}}$}

Set $u=x^2$; $R_u = 3$.  So $|x|<\sqrt{3}$.

\textit{Endpoints} $x=\pm\sqrt{3}$: $\sum\dfrac{1}{n\sqrt{n}}=\sum n^{-3/2}$ converges ($p=3/2$).

\textbf{Domain}: $[-\sqrt{3},\,\sqrt{3}]$.

\subsection*{(h) \hard\ $\displaystyle\sum_{n=0}^{\infty}\frac{(-1)^n \ln(n+2)(2x+3)^n}{(n+1)(2n)!}$}

The coefficients involve $(2n)!$ in the denominator, so the Ratio Test gives $R=+\infty$ for the variable $u=2x+3$.  Since this holds for all $u\in\mathbb{R}$, the series converges for all $x$.

\textbf{Domain}: $(-\infty,+\infty)$.

\subsection*{(i) \hard\ $\displaystyle\sum_{n=1}^{\infty}\frac{1\cdot3\cdots(2n-1)}{2\cdot4\cdots(2n)}\,x^n$}

Ratio Test: $\dfrac{a_{n+1}}{a_n}=\dfrac{2n+1}{2n+2}\to 1$.  So $R=1$.

\textit{Endpoints}: by Wallis, $a_n\sim 1/\sqrt{\pi n}$, so $\sum a_n$ diverges at both $x=\pm 1$.

\textbf{Domain}: $(-1, 1)$.

\subsection*{(j) \hard\ $\displaystyle\sum_{n=0}^{\infty}\frac{(2n)!\,x^n}{4^n(n!)^2}$}

Same coefficients as above (recognising the central binomial).  Ratio Test gives $R=1$; endpoints diverge.

\textbf{Domain}: $(-1, 1)$.

\subsection*{(k) \medium\ $\displaystyle\sum_{n=1}^{\infty}\frac{(3n-1)(-1)^n x^n}{n^3}$}

$R = \lim\dfrac{|a_n|}{|a_{n+1}|}=\lim\dfrac{(3n-1)}{n^3}\cdot\dfrac{(n+1)^3}{3n+2}=1$.

\textit{Endpoints}:
\begin{itemize}
  \item $x=1$: $\sum(3n-1)(-1)^n/n^3$.  Since $(3n-1)/n^3\searrow 0$, \textbf{Leibniz} applies; converges.
  \item $x=-1$: $\sum(3n-1)/n^3\sim\sum 3/n^2$; converges ($p=2$).
\end{itemize}
\textbf{Domain}: $[-1, 1]$.

% -------------------------------------------------------
\section*{Exercise 2 — Taylor Series around a Point $a$}

\subsection*{(a) $e^x$ around $a=0$}
$e^x = \displaystyle\sum_{n=0}^{\infty}\frac{x^n}{n!}$, $R=\infty$.

\subsection*{(b) $\sin x$ around $a=0$}
$\sin x = \displaystyle\sum_{n=0}^{\infty}\frac{(-1)^n x^{2n+1}}{(2n+1)!}$, $R=\infty$.

\subsection*{(c) $\dfrac{1}{1-x}$ around $a=0$}
$\dfrac{1}{1-x} = \displaystyle\sum_{n=0}^{\infty} x^n$, $|x|<1$.

\subsection*{(d) \medium\ $\dfrac{1}{x-5}$ around $a=1$}

Write $x-5=(x-1)-4$, so:
\[
  \frac{1}{x-5}=\frac{1}{(x-1)-4}=-\frac{1}{4}\cdot\frac{1}{1-\frac{x-1}{4}}
  =-\sum_{n=0}^{\infty}\frac{(x-1)^n}{4^{n+1}}, \quad |x-1|<4.
\]
\textbf{Domain}: $(-3, 5)$.

\subsection*{(e) \medium\ $\sin x$ around $a=\pi/2$}

Let $t=x-\pi/2$.  Then $\sin x=\sin(t+\pi/2)=\cos t=\sum_{n=0}^{\infty}\dfrac{(-1)^n t^{2n}}{(2n)!}$:
\[
  \sin x = \sum_{n=0}^{\infty}\frac{(-1)^n(x-\pi/2)^{2n}}{(2n)!}, \quad R=\infty.
\]

\subsection*{(f) \easy\ $\ln x$ around $a=1$}

$\ln x = \ln(1+(x-1)) = \displaystyle\sum_{n=1}^{\infty}\frac{(-1)^{n+1}(x-1)^n}{n}$, $\;0 < x \leq 2$.

\subsection*{(g) \hard\ $\dfrac{1}{(x-5)^3}$ around $a=1$}

From part (d): $\dfrac{1}{x-5}=-\sum_{n=0}^{\infty}\dfrac{(x-1)^n}{4^{n+1}}$.

Differentiate twice with respect to $x$:
\[
  \frac{1}{(x-5)^2} = \sum_{n=1}^{\infty}\frac{n(x-1)^{n-1}}{4^{n+1}},\qquad
  \frac{2}{(x-5)^3} = \sum_{n=2}^{\infty}\frac{n(n-1)(x-1)^{n-2}}{4^{n+1}}.
\]
Dividing by 2 and re-indexing ($m=n-2$):
\[
  \frac{1}{(x-5)^3} = \sum_{m=0}^{\infty}\frac{(m+2)(m+1)}{2\cdot 4^{m+3}}(x-1)^m, \quad |x-1|<4.
\]

\subsection*{(h) \medium\ $\dfrac{1}{(x-2)^2}$ around $a=-3$}

Write $x-2=(x+3)-5$, so $\dfrac{1}{x-2}=-\dfrac{1}{5}\cdot\dfrac{1}{1-(x+3)/5}=-\sum_{n=0}^{\infty}\dfrac{(x+3)^n}{5^{n+1}}$.

Differentiate with respect to $x$:
\[
  \frac{1}{(x-2)^2} = \sum_{n=1}^{\infty}\frac{n(x+3)^{n-1}}{5^{n+1}}, \quad |x+3|<5.
\]
\textbf{Domain}: $(-8, 2)$.

% -------------------------------------------------------
\section*{Exercise 3 — Maclaurin Series (Substitution)}

\subsection*{(a) \easy\ $e^{-x^2}$}
Substitute $u=-x^2$ into $e^u=\sum u^n/n!$:
\[
  e^{-x^2} = \sum_{n=0}^{\infty}\frac{(-1)^n x^{2n}}{n!}, \quad \forall x\in\mathbb{R}.
\]

\subsection*{(b) \easy\ $\sin(x^3)$}
Substitute $u=x^3$ into $\sin u=\sum(-1)^n u^{2n+1}/(2n+1)!$:
\[
  \sin(x^3) = \sum_{n=0}^{\infty}\frac{(-1)^n x^{6n+3}}{(2n+1)!}, \quad \forall x\in\mathbb{R}.
\]

\subsection*{(c) \medium\ $\dfrac{1}{x^2+4x+5}$}

Complete the square: $x^2+4x+5=(x+2)^2+1$.  Let $u=-(x+2)^2$:
\[
  \frac{1}{(x+2)^2+1} = \frac{1}{1-u}\Big|_{u=-(x+2)^2}
  = \sum_{n=0}^{\infty}(-1)^n(x+2)^{2n}.
\]
This is a Taylor series centered at $a=-2$, not $a=0$.  For a Maclaurin series ($a=0$), write $v=(4x+x^2)/5$ so $x^2+4x+5=5(1+v)$:
\[
  \frac{1}{5}\cdot\frac{1}{1+v} = \frac{1}{5}\sum_{n=0}^{\infty}(-1)^n v^n
  = \frac{1}{5}\sum_{n=0}^{\infty}(-1)^n\!\left(\frac{x^2+4x}{5}\right)^{\!n},
  \quad |x^2+4x|<5.
\]

\subsection*{(d) \easy\ $x e^x$}
Multiply $e^x=\sum x^n/n!$ by $x$:
\[
  xe^x = \sum_{n=0}^{\infty}\frac{x^{n+1}}{n!} = \sum_{n=1}^{\infty}\frac{x^n}{(n-1)!}, \quad \forall x.
\]

\subsection*{(e) \easy\ $\ln(1+x^2)$}
Substitute $u=x^2$ into $\ln(1+u)=\sum_{n=1}^{\infty}(-1)^{n+1}u^n/n$:
\[
  \ln(1+x^2) = \sum_{n=1}^{\infty}\frac{(-1)^{n+1}x^{2n}}{n}, \quad |x|\leq 1.
\]

\subsection*{(f) \hard\ $\arctan\!\left(\dfrac{1+x^2}{1-x^2}\right)$}

Use the tangent addition identity: $\dfrac{1+x^2}{1-x^2} = \tan\!\left(\dfrac{\pi}{4}+\arctan(x^2)\right)$ (verify by $\tan A+\tan B$ formula with $A=\pi/4$, $\tan A=1$, $\tan B=x^2$, $\tan(A+B)=(1+x^2)/(1-x^2)$).

Therefore $\arctan\!\left(\dfrac{1+x^2}{1-x^2}\right) = \dfrac{\pi}{4}+\arctan(x^2)$.

Using $\arctan u=\sum_{n=0}^{\infty}\dfrac{(-1)^n u^{2n+1}}{2n+1}$ and substituting $u=x^2$:
\[
  f(x) = \frac{\pi}{4} + \sum_{n=0}^{\infty}\frac{(-1)^n x^{4n+2}}{2n+1}, \quad |x|<1.
\]

% -------------------------------------------------------
\section*{Exercise 4 — High-Order Derivatives via Taylor Coefficients}

\subsection*{(a) \medium\ $f^{(17)}(1)$ for $f(x)=\dfrac{1}{x-5}$}

From Exercise 2(d): $\dfrac{1}{x-5}=-\sum_{n=0}^{\infty}\dfrac{(x-1)^n}{4^{n+1}}$.

The Taylor coefficient of $(x-1)^{17}$ is $-1/4^{18}$.  Since $c_{17}=f^{(17)}(1)/17!$:
\[
  f^{(17)}(1) = 17!\cdot\left(-\frac{1}{4^{18}}\right) = \boxed{-\frac{17!}{4^{18}}}.
\]

\subsection*{(b) \medium\ $\left.\dfrac{d^{10}}{dx^{10}}\!\left(\dfrac{1}{1-x}\right)\right|_{x=0}$}

The Maclaurin series is $\sum_{n=0}^{\infty}x^n$.  The coefficient of $x^{10}$ is $1$.
Since $c_{10} = f^{(10)}(0)/10!$:
\[
  f^{(10)}(0) = 10!\cdot 1 = \boxed{10!} = 3\,628\,800.
\]

\subsection*{(c) \hard\ Evaluating $\sum_{n=0}^{\infty}\frac{f^{(n)}(a)}{n!}(1-a)^n$}

This expression is precisely the Taylor series of $f$ centered at $a$, evaluated at $x=1$:
\[
  \sum_{n=0}^{\infty}\frac{f^{(n)}(a)}{n!}(1-a)^n = f(1),
\]
provided $|1-a| < R_a$ (the radius of convergence of the Taylor series at $a$).  Since by hypothesis $R_a > 1$, we have $|1-a|=|1-a|\leq 1 < R_a$, so the series converges to $f(1)$. $\square$

\bigskip
\hrule
\bigskip
\textit{Key takeaways}: The radius $R$ is found by Ratio/Root test on coefficients.  Endpoint analysis requires separate care.  Substitution is the fastest way to build Maclaurin series for composite functions.  Taylor coefficients directly give all derivatives at the center: $f^{(n)}(a)=n!\cdot c_n$.

\end{document}
